\section{竞争格局与差异化}
\label{sec:competition}

\subsection{三种角色，一个空白}

市场上与 Vulca \keyword{"看起来类似"}的厂商可以归为两类，
但它们\keyword{都不与 Vulca 直接竞争}，原因是各自选择的技术哲学和商业模型不同。
这个二维坐标把整个战场讲清楚：

\begin{center}
\small
\begin{tabular}{@{}p{3.2cm}p{4.8cm}p{4.8cm}@{}}
\toprule
 & \heiti\color{vulcaPrimary}\quad 封闭（整合式） & \heiti\color{vulcaPrimary}\quad 开放（可组合） \\
\midrule
\heiti 代表 & \textbf{OpenAI Images / Adobe Firefly / Runway} & \textbf{HuggingFace / SAM / ControlNet / Diffusers} \\
\heiti 产品定位 & 自带"大脑"决策，全流程封闭 & 提供"零件"，不管工作流 \\
\heiti 面向 & 终端消费者 / 企业一站式 & 开发者自己拼装 \\
\heiti 对 agent & \keyword{agent 的竞争者}（自身就是 agent） & \keyword{agent 的原子能力}（不是工作流） \\
\midrule
\heiti\color{vulcaPrimary}\quad Vulca 的位置 & \multicolumn{2}{l}{\keyword{开放 + 整合式：给 agent 用的"手"；不自己决策，但提供完整工作流。}} \\
\bottomrule
\end{tabular}
\end{center}

\subsection{关键竞品拆解}

\subsubsection*{"大脑派"：OpenAI / Adobe Firefly / Runway}
\begin{itemize}
  \item \keyword{定位}：端到端产品，自己做 UI、自己做规划、自己做生成；
  \item \keyword{对 Vulca 的关系}：不是竞品，是\keyword{共存关系}——它们面向不用 agent 的终端用户，
  而 Vulca 面向\keyword{已经用 Claude Code / Cursor / 自建 agent 的专业人群}；
  \item \keyword{在 agent 生态里}：它们的 API 可以作为 Vulca 的\keyword{后端 provider 之一}（OpenAI provider 已支持），
  但不提供 MCP 原生工具化；
  \item \keyword{国内市场}：均不提供本地化部署，数据出境风险高，政府与大机构采购受限。
\end{itemize}

\subsubsection*{"零件派"：HuggingFace / SAM / ControlNet / Diffusers}
\begin{itemize}
  \item \keyword{定位}：提供模型和基础算法，开发者自己写脚本拼装；
  \item \keyword{对 Vulca 的关系}：\keyword{上游依赖}而非竞品——Vulca 用 EVF-SAM、SDXL、FLUX 等作为 provider 后端；
  \item \keyword{差距}：它们不提供\keyword{agent-native 调用}（MCP）、\keyword{工作流封装}（skill）、
  \keyword{文化评估框架}（L1–L5 YAML），三者都是 Vulca 的增量。
\end{itemize}

\subsubsection*{"agent-native 图像"：尚无直接竞品}
据公开信息梳理，截至 2026-04-23：

\begin{itemize}
  \item 没有公司以"agent 的手"作为产品定位；
  \item 现有 MCP 生态中的图像工具以\keyword{单一能力包装}（如"把文生图暴露成 MCP"）为主，
  没有提供\keyword{完整工作流链路}；
  \item 没有将\keyword{文化评估}作为原生能力嵌入 SDK 的开源项目；
  \item 国内尚无团队在 Claude Code / MCP 插件生态位置上占有显著存在。
\end{itemize}

\subsubsection*{邻近项目对比（海内外共 5 家）}

为避免"尚无直接竞品"被误读为调研不足，我们按\keyword{五个维度}对比常被归为同赛道的邻近项目。
结论是：每家都与 Vulca 的"\textbf{agent-native 像素执行层}"存在一个以上错位，\keyword{五个维度同时具备的只有 Vulca}。

\begin{center}
\scriptsize
\begin{tabularx}{\linewidth}{@{}>{\tabRR}p{2.1cm}|>{\tabRR}p{1.4cm}|>{\tabRR}p{1.6cm}|>{\tabRR}p{1.6cm}|>{\tabRR}p{1.4cm}|>{\tabRR}X@{}}
\toprule
\rowcolor{tableHead}
\heiti 项目 & \heiti MCP 原生 & \heiti 工作流封装 & \heiti 文化评估 & \heiti 本地部署 & \heiti 商业定位 \\
\midrule
\keyword{ComfyUI} & 否（社区插件零散） & 节点图可复用 & 否 & 是 & 开源节点引擎 / 社区 \\
\keyword{Replicate} & 否（REST API） & 否（单一模型封装） & 否 & 否 & 云端模型托管 \\
\keyword{Stability AI SDK} & 否 & 否（模型调用） & 否 & 部分 & 基础模型厂商 \\
\keyword{即梦 / 文心一格 / 商汤秒画} & 否 & 否（应用层产品） & 否 & 否 & C 端 / B 端生成工具 \\
\keyword{Adobe Firefly Services} & 否 & 封闭 workflow & 否 & 否 & 企业端封闭 SaaS \\
\midrule
\rowcolor{vulcaSoft}
\keyword{\textbf{Vulca}} & \textbf{是（21 工具）} & \textbf{是（3 shipped skill）} & \textbf{是（13 种 L1–L5）} & \textbf{是（MPS 已验证）} & \textbf{开源 SDK + 托管 + On-Prem} \\
\bottomrule
\end{tabularx}
\end{center}

\paragraph{补充说明}
\begin{itemize}
  \item \keyword{ComfyUI} 与 Vulca \textbf{非竞品}，是 Vulca 的 \keyword{provider 后端之一}（local 路径核心组件）；
  \item \keyword{Replicate / Stability} 提供模型托管或 API，与 Vulca 的\keyword{工作流编排}层不重叠；
  \item \keyword{国内 C 端 / 大厂应用}（即梦、文心一格、商汤秒画等）与 Vulca 在\keyword{"应用层 vs SDK/agent 插件层"}错位——它们的用户是终端创作者，Vulca 的"用户"是 agent；
  \item \keyword{Adobe Firefly Services} 走封闭企业 SaaS 路径，不开放 MCP，也不支持本地部署；
  \item \keyword{其它 MCP 生态里的图像工具}目前以\keyword{单一能力包装}（如"把文生图暴露成 MCP"）为主，\keyword{没有一家}提供完整的"拆图 $\to$ 编辑 $\to$ 评估 $\to$ 合成"链路加文化纵深。
\end{itemize}

\subsection{Vulca 的四项差异化}

\begin{enumerate}
  \item \keyword{agent-native 协议先行}：21 个 MCP 工具原生支持 Claude Code 插件生态，
  不是"另做一个 UI 加 MCP 包装"；
  \item \keyword{文化纵深}：13 种艺术/设计传统的 L1–L5 评价体系，原生覆盖中国工笔/写意等东方语境；
  \item \keyword{自主可控}：完全本地化链路 + Apache 2.0，
  数据不出网、算法可审计、硬件可降级；
  \item \keyword{工程质量信号}：高覆盖率测试套件、3 个已上线 skill、Phase 0/1 两次架构重构零回归，
  不是\keyword{"发布即废弃"}的 demo 仓库。
\end{enumerate}

\subsection{壁垒强度与持续性}

\begin{itemize}
  \item \keyword{3–6 个月可复制}的部分：MCP 工具接口本身（但需要\keyword{生态位}与\keyword{信任}积累，信任无法靠代码堆出）；
  \item \keyword{6–12 个月可复制}的部分：本地化链路工程（特别是 MPS / 国产算力适配），\keyword{单一文化传统的 YAML 定义}（配一位文化顾问 + 学术资料可在一季度内完成）；
  \item \keyword{12 个月以上需专门投入}的部分：\keyword{13 种文化传统跨学科覆盖} + \keyword{EMNLP / ACM MM 等学术合作关系网} + \keyword{开源贡献者社区信任}——
  这类资产不是"不可复制"，而是\keyword{需要持续时间 + 专门学术投入}才能积累。竞品雇 3 位文化顾问一个季度可以追赶一种传统；追赶 13 种 + 学术认同需要更长周期。
\end{itemize}
